% Замените поясняющий текст собственными результатами.
% В шаблоне нет готового решения и обязательных внешних изображений.
\reportelement{Введение}

Целью работы является подключение к приложению PostgreSQL, обучение хранения данных вне контейнера и разбирание с сетями Docker.; персональные параметры указаны в таблице~\ref{tab:parameters}.

\begin{table}[H]
  \caption{Персональные параметры}\label{tab:parameters}
  \centering
  \begin{tabularx}{\textwidth}{|p{0.46\textwidth}|X|}
    \hline
    Параметр & Значение\\ \hline
    Номер ИСУ & \StudentISU\\ \hline
    Персональный порт (PostgreSQL) & 5489\\ \hline
    Персональный порт (приложение-ассистент) & 8059\\ \hline
    Персональный префикс & matveev\\ \hline
    Секретный ключ & ------\\ \hline
    Имя базы данных & matveev59\\ \hline
  \end{tabularx}
\end{table}

Результат хранится в ветке \texttt{lab-\LabNumber} форка:
\url{\RepositoryURL}. 

\chapter{Ход выполнения}
\section{Предварительные требования}

Создал именованный volume для БД, проверил создание.  Запустил PostgreSQL с volume, проверил статус. Запустил ассистент и подключил к БД на рисунке~\ref{fig:1}.

\reportscreenshot{Screenshot_1.png}{Запуск}{fig:1}

Проверка работоспсобности изображена на рисунке~\ref{fig:2}.

\reportscreenshot{Screenshot_2.png}{Проверка}{fig:2}



\section{Тест персистентности: пересоздание контейнера}

Удалил контейнер с ассистентом. Пересоздал контейнер ассистента. Проверил, что данные сохранились на рисунке~\ref{fig:3}.

\reportscreenshot{Screenshot_3.png}{Выполнение команд}{fig:3}

Артефакт персистентности изображен на рисунке~\ref{fig:4}.

\reportscreenshot{Screenshot_4.png}{Артефакт}{fig:4}

\section{Bind mount для конфигурации}

В папке config/ создал файл app.conf и создал контейнер с bind mount на рисунке~\ref{fig:5}.

\reportscreenshot{Screenshot_5.png}{Создание}{fig:5}

Проверил доступ к конфигурации на рисунке~\ref{fig:6}.

\reportscreenshot{Screenshot_6.png}{Проверка}{fig:6}

\section{Пользовательские сети для изоляции}

Создал пользовательскую сеть и пересоздал контейнеры в пользовательской сети на рисунке~\ref{fig:7}.

\reportscreenshot{Screenshot_7.png}{Создание}{fig:7}

Проверил связность на рисунке~\ref{fig:8}.

\reportscreenshot{Screenshot_8.png}{Проверка}{fig:8}

Исследовал сетевую изоляцию на рисунке~\ref{fig:9} с запуском временного контейнера вне и внутри сети соответственно.

\reportscreenshot{Screenshot_9.png}{Исследование сети}{fig:9}

\section{Финальный артефакт}

Финальная проверка инфраструктуры изображена на рисунке~\ref{fig:10}.

\reportscreenshot{Screenshot_10.png}{Проверка}{fig:10}

\reportelement{Заключение}

В ходе работы к приложению-ассистенту была подключена база данных PostgreSQL, а данные вынесены за пределы контейнера в именованный том \texttt{matveev-pgdata}. Тест персистентности подтвердил, что заметки сохраняются после удаления и повторного создания контейнера ассистента. С помощью bind mount конфигурация подключена в режиме \texttt{read-only} и меняется без пересборки образа. Перенос контейнеров в пользовательскую сеть \texttt{matveev-backend} позволил обращаться к базе по имени контейнера, а проверка \texttt{ping} показала изоляцию: контейнер вне сети базу не видит.

% Оставьте только использованные источники и цитируйте их в тексте.
\begin{thebibliography}{9}
  \bibitem{docker-docs}
  Docker Documentation: документация Docker.
  URL: \url{https://docs.docker.com/}
  (дата обращения: укажите свою дату).
\end{thebibliography}

